% part: intuitionistic-logic
% chapter: semantics
% section: topological-semantics

\documentclass[../../../include/open-logic-section]{subfiles}

\begin{document}

\olfileid{int}{sem}{top}

\olsection{స్థలవిజ్ఞాన అర్థవిజ్ఞానం}

అంతఃప్రజ్ఞావాద తర్కానికి అర్థవిజ్ఞానం ఇవ్వడానికి
మరో మార్గం స్థలవిన్యాసం అనే గణిత భావనను వాడటం.

\begin{defn}
  $X$ ఒక సమితి అనుకుందాం. $X$పై ఒక
  \emph{స్థలవిన్యాసం (టోపాలజీ)} అనేది కింది
  ధర్మాలు గల సమితి $\Top{O}
  \subseteq \Pow{X}$. $\Top{O}$ యొక్క
  \tetoken{మూలకాలను}{!!{element}s}
  స్థలవిన్యాసపు \emph{వివృత సమితులు} అంటాం.
  $X$, $\Top{O}$ జతను \emph{స్థలవిజ్ఞాన
  అంతరాళం} అంటాం.
  \begin{enumerate}
  \item శూన్య సమితి, మొత్తం అంతరాళం వివృతమే:
    $\emptyset$, $X \in
    \Top{O}$.
  \item వివృత సమితులు పరిమిత ఛేదనల కింద
    మూసుకుపోయి ఉంటాయి: $U$, $V \in
    \Top{O}$ అయితే $U \cap V \in \Top{O}$.
  \item వివృత సమితులు యథేచ్ఛ సమ్మేళనాల కింద
    మూసుకుపోయి ఉంటాయి: అన్ని $i \in I$లకు
    $U_i \in
    \Top{O}$ అయితే $\bigcup \Setabs{U_i}{i \in
      I} \in \Top{O}$.
  \end{enumerate}
\end{defn}

వివృత సమితుల సముదాయాన్ని సందర్భం నుంచి
తెలుసుకోవచ్చంటే, స్థలవిన్యాసానికి $X$ అనే
సంకేతాన్నే వాడవచ్చు. అయినా $X$కు వివృత
సమితుల సముదాయం ఇచ్చిన తరువాతనే దాన్ని
స్థలవిన్యాసంగా పిలవవచ్చని గమనించండి.

\begin{defn}
  అంతఃప్రజ్ఞావాద ప్రతిజ్ఞావాక్య తర్కపు
  \emph{స్థలవిజ్ఞాన నమూనా} ఒక త్రయం
  $\mModel{X} = \tuple{X, \Top{O}, V}$;
  ఇందులో $\Top{O}$ అనేది~$X$పై స్థలవిన్యాసం,
  $V$ అనేది ప్రతి ప్రతిజ్ఞావాక్య చరానికి
  $\Top{O}$లోని ఒక వివృత సమితిని కేటాయించే
  ప్రమేయకం.

  స్థలవిజ్ఞాన నమూనా~$\mModel{X}$ ఇచ్చినప్పుడు,
  $\Prop{X}{!A}$ను ఇలా ఆగమన పద్ధతిలో
  నిర్వచించవచ్చు:
  \begin{enumerate}
  \item $\Prop{X}{\lfalse} = \emptyset$
  \item $\Prop{X}{p} = V(p)$
  \item $\Prop{X}{!A \land !B} = \Prop{X}{!A} \cap \Prop{X}{!B}$
  \item $\Prop{X}{!A \lor !B} = \Prop{X}{!A} \cup \Prop{X}{!B}$
  \item $\Prop{X}{!A \lif !B} = \Interior{(X \setminus \Prop{X}{!A}) \cup
    \Prop{X}{!B}}$
  \end{enumerate}
  ఇక్కడ $\Interior{V}$ అనేది సమితి $V \subseteq X$ను
  దాని \emph{అంతర్భాగం}కు పంపే ప్రమేయకం;
  అంటే ఆ సమితిలో ఉన్న అన్ని వివృత సమితుల
  సమ్మేళనం. మరో మాటలో,
  \[
  \Interior{V} = \bigcup \Setabs{U}{U \subseteq V \text{ మరియు } U \in \Top{O}}.
  \]
\end{defn}

ఏ సమితి అంతర్భాగమైనా ఎల్లప్పుడూ వివృతమే;
ఎందుకంటే అది వివృత సమితుల సమ్మేళనం. అందువల్ల
$\Prop{X}{!A}$ ఎల్లప్పుడూ వివృత సమితి.

స్థలవిజ్ఞాన అర్థవిజ్ఞానం చాలా అమూర్తమైనదైనా,
దానికి ప్రేరణ ఇచ్చే విధంగా ఆలోచించవచ్చు.
$X$ యొక్క \tetoken{మూలకాలు}{!!{element}s},
లేదా ``బిందువులు'', వాక్యాలను మూల్యాంకనం
చేయగల స్థానాలనుకుందాం. $!A$ సత్యమయ్యే
అన్ని బిందువుల సమితి,~$!A$ వ్యక్తీకరించే
ప్రతిపాదన. ప్రతి బిందు సమితీ సాధ్యమైన
ప్రతిపాదన కాదు; $\Top{O}$ యొక్క
\tetoken{మూలకాలు}{!!{element}s} మాత్రమే అలాంటివి.
ప్రతి~$X$కు $\Prop{X}{!A}
\subseteq \Prop{X}{!B}$ అయితే, అప్పుడే $!A \Entails !B$;
అంటే~$!A$ సత్యమయ్యే ప్రతి బిందువు వద్దా
$!B$ సత్యం. అసంభవ వాక్యం~$\lfalse$
ఎప్పుడూ సత్యం కాదు; కాబట్టి
$\Prop{X}{\lfalse} = \emptyset$.

అంతఃప్రజ్ఞావాదంలో చెల్లుబాటయ్యే సూత్రాలు
నిలవాలంటే, $!B \land !C$, $!B \lor !C$,
$!B \lif !C$ వ్యక్తీకరించే ప్రతిపాదనలు,
$!B$,~$!C$ ప్రతిపాదనలతో ఎలా సంబంధించాలి?
అంటే $!A \Proves
!B$ అయితే, అప్పుడే $\Prop{X}{!A} \subset \Prop{X}{!B}$
కావాలంటే? ఏ $!A$కైనా $\lfalse
\Proves !A$ కావాలి; ప్రతి $U$కు $\emptyset
\subseteq U$ కాబట్టి అది నెరవేరుతుంది.
$!B \land !C \Proves !B$ కాబట్టి
$\Prop{X}{!B \land !C} \subseteq \Prop{X}{!B}$
కావాలి; అలాగే $\Prop{X}{!B \land !C}
\subseteq \Prop{X}{!C}$ కావాలి. $W \subseteq U$,
$W \subseteq V$ రెండూ నెరవేర్చే అతి పెద్ద
సమితి $U \cap V$. మరోవైపు
$!B \Proves !B \lor !C$, $!C \Proves !B \lor !C$;
కాబట్టి $\Prop{X}{!B} \subseteq \Prop{X}{!B \lor !C}$,
$\Prop{X}{!C} \subseteq \Prop{X}{!B \lor !C}$
కావాలి. $U \subseteq W$, $V \subseteq W$
అయ్యే అతి చిన్న సమితి~$W$ అనేది $U \cup V$.

$\lif$ నిర్వచనం కొంత క్లిష్టం: $!A \lif !B$
అనేది $!A$తో కలిసినప్పుడు $!B$ను
అనుగమింపజేసే అతి బలహీన ప్రతిపాదన.
$(!A \lif !B)
\land !A \Proves !B $ కాబట్టి $!A
\lif !B$ను $!A$తో కలిపితే~$!B$ వస్తుందని
స్పష్టం. కాబట్టి $\Prop{X}{!A \lif !B} \cap \Prop{X}{!A}
\subset \Prop{X}{!B}$ అయ్యే అతి పెద్ద వివృత
సమితి $\Prop{X}{!A \lif !B}$ కావాలి; ఇదే
మన నిర్వచనానికి దారితీస్తుంది.

\end{document}
